\chapter{Introduction and literature context}
An unbounded isotropic solid supports longitudinal and transverse waves. In a
plate, boundary reflection and mode conversion produce discrete guided branches.
Lamb's plate solution established their characteristic equations
\cite{lamb1917,viktorov1967}; modern treatments emphasize dispersion, multimode
propagation, and defect sensitivity \cite{graff1975,rose2014}.

Laser ultrasonics provides noncontact generation. Below damage threshold,
absorbed light causes thermal expansion; at higher fluence ablation adds recoil.
Only the nondestructive thermoelastic regime is treated here \cite{scruby1990}.
Arias and Achenbach's line-source theory includes finite width, duration, optical
penetration, and diffusion \cite{arias2003}. ZGV resonances were generated and
detected optically by Prada et al. \cite{prada2005} and used for elastic
characterization \cite{clorennec2007}. A moving continuous-wave laser can select
branches through phase-velocity matching \cite{li2018}.

A frequency-domain solution assumes a source acting indefinitely. For activation
at $t=0$, four notions differ: first arrival; decay of the homogeneous solution;
convergence of a measurement to its forced periodic limit; and resonance-
controlled buildup near a cutoff or ZGV. This thesis derives these distinctions
for switched sinusoids, bursts, scans, and arrays. Its target observables are a
detector-aware settling time, a defined burst length, and modal selectivity from
physical thermoelastic overlaps. Current results reach verified dispersion only;
later levels are identified explicitly in \cref{ch:results}.
